The results distinguish agreement with the crowd-vote distribution from disagreement among accepted category predictions. Distributional supervision improves the former across the tested source-dataset configurations, but its selective advantage does not persist across both backbones. Fixed-image comparisons further show that the predictor ordering depends on the accepted set. These findings answer RQ1 by qualifying the primary Swin-T advantage: distributional distance and selective risk describe different aspects of performance and cannot be treated as interchangeable evidence of improved recognition.

This distinction develops the evaluation question raised by prior distribution-based FER studies. FER+ and crowdsourced affective supervision established the value of retaining human response distributions rather than treating all disagreement as labeling error \cite{barsoum2016ferplus,washington2021crowdsourcing}. The present controls narrow the interpretation of a hard-versus-soft comparison by preserving dominant target probabilities and tied maxima. They show a smaller selective advantage than the hard-target comparison while retaining a consistent distributional advantage. That pattern supports evaluating both outcomes, but the controls still differ in entropy and aggregate class proportions. The experiment compares these specified target representations rather than isolating the contribution of minority-category identity.

The primary selective effect also warrants attention to scale. A difference of 0.38 pp across 2,620 accepted images corresponds to about 9.9 fewer expected disagreements with a vote drawn from the observed annotations, not 9.9 corrected images. Losses are fractional and the accepted identities differ. For RQ2, the fixed-image analysis explains why each model's own-ranking risk alone is insufficient: prediction quality on a specified accepted set and the composition of that set contribute jointly to the reported outcome. Reference-specific F1 and unsupported-category accounting add two further requirements for interpretation---the label reference and the complete annotation vocabulary must be explicit.

The confidence-score and transfer checks delimit the scope of the selective result for RQ3. Mean differences remain backbone-dependent under the three probability-only scores, and some individual seed comparisons change direction. This extends the connection between annotation disagreement and selective evaluation studied in Crowd-Calibrator \cite{khurana2024crowd}, while leaving learned rejection methods such as SelectiveNet outside the implemented comparison \cite{geifman2019selectivenet}. On RAVDESS, the ordering also depends on the fitted checkpoint and on whether calm and neutral responses are merged. The external diagnostics identify substantial excess disagreement and neutral-dominated predictions, but do not assign a causal explanation to capture conditions, population or temporal aggregation.

Finally, the calibration-transfer results distinguish a nominal risk target from achieved risk. Greater coverage at an empirical threshold is not an equal-risk advantage when the transferred risks differ. The more conservative comparator makes its coverage cost explicit, without verifying that its assumptions hold in deployment. This is consistent with the distinction between confidence scores and calibrated probabilities \cite{guo2017calibration}. Together, the findings motivate a coherent evaluation practice: compare matched supervision targets, report distributional and selective outcomes separately, examine identical accepted sets, specify the annotation reference and vocabulary, and present achieved risk alongside coverage. The conditional intervals quantify sampling sensitivity for the retained fits and observed annotations; they do not establish a general ordering over future training seeds, raters or populations.
